\section{Discussion}

Nghiên cứu này sử dụng các phương pháp phân tích chuỗi thời gian phi tuyến bổ sung cho nhau để đặc trưng động lực học PPG trên các cửa sổ ngắn khi chuyển từ Awake sang Drowsy. Kết quả cho thấy, ở cả hai trạng thái, tín hiệu PPG chứa tổ chức theo thời gian vượt quá những gì một noisy pseudoperiodic process có thể tái tạo, nhưng điều này không đồng nghĩa với bằng chứng của deterministic chaos. Đồng thời, Drowsiness đi kèm với giảm short-term forecastability, giảm tổ chức recurrence theo đường chéo và giảm local trajectory divergence, trong khi các chỉ số laminar recurrence có xu hướng tăng. Nhìn chung, các phát hiện này cho thấy quá trình Awake--Drowsy liên quan đến sự tái tổ chức đồng thời của nhiều thuộc tính động lực học khác nhau của PPG, thay vì một thay đổi đơn giản theo trục more chaos'' hay less chaos''~\cite{porta2007_prediction,porta2015,calderonjuarez2023,hu2009}.

\subsection{Dynamical Changes in PPG Across the Awake--Drowsy Transition}

Kết quả prediction cho thấy Drowsiness đi kèm với prediction correlation thấp hơn và normalized prediction error cao hơn, tương ứng với reduced finite-horizon forecastability. Trong local prediction, các trạng thái gần nhau trong reconstructed state space được kỳ vọng tạo ra những future trajectories tương tự nếu tổ chức theo thời gian đủ ổn định. Vì vậy, kết quả hiện tại cho thấy diễn tiến của PPG trong trạng thái Drowsy khó được dự báo từ các trạng thái gần trong quá khứ hơn so với trạng thái Awake. Cách diễn giải này phù hợp với Porta et al., trong đó local prediction được sử dụng để đặc trưng short-term cardiovascular predictability và được phân biệt rõ với nonlinearity hay deterministic chaos \cite{porta2007_prediction}. Faes et al.\ cũng cho thấy predictability của cardiovascular dynamics có thể phản ánh cả self-dynamics lẫn tương tác với các hệ sinh lý khác \cite{faes2016}. Do đó, reduced forecastability trong Drowsiness có thể phản ánh sự thay đổi của tổ chức theo thời gian trong một regulatory configuration đang biến đổi, chứ không nhất thiết biểu thị greater complexity hoặc more chaos.

RQA cung cấp một góc nhìn hình học bổ sung. Drowsiness đi kèm với DET giảm và Lmean có xu hướng giảm, trong khi LAM và TT có xu hướng tăng. Diagonal recurrence structures phản ánh sự lặp lại của diễn tiến quỹ đạo theo thời gian, còn vertical structures phản ánh local persistence hoặc trapping trong các vùng của reconstructed state space. Các nghiên cứu RQA liên quan đến sleep trước đây đã cho thấy cardiovascular recurrence organization thay đổi theo trạng thái sinh lý \cite{rolink2015}, trong khi Martín-González et al.\ cho thấy các chỉ số theo đường chéo và theo chiều dọc có thể phản ứng khác nhau và chứa các thông tin động lực học khác nhau \cite{martingonzalez2018}. Pattern hiện tại vì vậy phù hợp hơn với một sự tái phân bố của recurrence geometry: tổ chức theo đường chéo suy yếu trong khi laminar persistence có xu hướng tăng. DET giảm không nên được diễn giải là giảm physical determinism, và LAM/TT tăng cũng không đồng nghĩa với nonlinearity lớn hơn \cite{calderonjuarez2023}.

LLE thấp hơn trong Drowsiness cho thấy local trajectory divergence giảm trong khoảng thời gian fit. Kết quả này không nên được hiểu là ``less chaos'', vì LLE phụ thuộc mạnh vào scale, embedding, data length, preprocessing và signal characteristics \cite{yeragani2004,hu2009}. Việc forecastability giảm đồng thời với LLE giảm cũng không phải mâu thuẫn: Simplex prediction và LLE đo các thuộc tính khác nhau của dynamics và có thể nhạy với các temporal scales khác nhau \cite{porta2015,faes2016,hu2009}. Vì vậy, toàn bộ kết quả phù hợp với một sự thay đổi đồng thời của nhiều dynamical properties trong Awake--Drowsy transition.

\subsection{Physiological Interpretation of the Awake--Drowsy Transition}

Các thay đổi trên nên được đặt trong bối cảnh wake-to-sleep transition, thay vì xem Drowsiness như một sleep stage ổn định. Sleep onset là một quá trình chuyển tiếp có sự tái tổ chức của central nervous, autonomic, respiratory, baroreflex và cardiovascular regulation \cite{dezambotti2018,tobaldini2013}. Shinar et al.\ cho thấy wake-like và sleep-like states có thể dao động trước khi stable sleep được thiết lập \cite{shinar2006}; Carrington et al.\ cho thấy repeated arousals có thể tạm thời khôi phục wake-like cardiovascular activity và làm gián đoạn quá trình settling \cite{carrington2005}; còn Tsai et al.\ cho thấy cardiovascular de-arousal có thể xảy ra trước cortical N1 onset \cite{tsai2019}. Vì vậy, Drowsiness trong dataset hiện tại phù hợp hơn với một transitional regulatory regime không ổn định hơn là stable NREM sleep.


Framework này cung cấp bối cảnh sinh lý hợp lý cho việc forecastability, recurrence geometry và local divergence cùng thay đổi. Một regulatory configuration đang thay đổi liên tục dưới ảnh hưởng của autonomic, respiratory, cardiovascular và cortical processes có thể làm biến đổi tổ chức theo thời gian của PPG. Tuy nhiên, literature hiện tại không trực tiếp chứng minh cơ chế gây ra $CC\downarrow$, $DET\downarrow$ hay $LLE\downarrow$; do đó đây vẫn là physiological interpretation, không phải causal explanation. Symbolic analysis cung cấp thêm bằng chứng hỗ trợ cho altered cardiac autonomic regulation. Pattern $0V\uparrow / 2V\downarrow$ tương tự direction được Porta et al.\ ghi nhận trong graded head-up tilt \cite{porta2007_symbolic}. Tuy nhiên, nhiều nghiên cứu sleep-onset lại báo cáo xu hướng tăng vagal contribution hoặc giảm LF/HF \cite{shinar2006,carrington2003,tsai2019,vaussenat2026}. Sự khác biệt này không nhất thiết là contradiction, vì sleep onset có thể chứa incomplete de-arousal, repeated arousals và sự chuyển đổi không đồng bộ giữa các hệ. Do đó, symbolic result chỉ nên được xem là phù hợp với altered autonomic modulation, không phải bằng chứng trực tiếp cho một sympathetic hoặc parasympathetic shift cụ thể.

\subsection{Ultra-Short PPG Dynamics and Practical Implications}

Direction của các effects chính về forecastability, DET và LLE được bảo toàn trên các cửa sổ 30, 60, 120 và 180~s, dù 30~s có uncertainty lớn hơn. Kết quả này phù hợp với Sviridova et al., người cho thấy minimum recording length của PPG RQA phụ thuộc mạnh vào từng metric và DET ổn định sớm hơn các diagonal-length measures \cite{sviridova2022}. Điều này cũng phù hợp với kết quả hiện tại khi DET bền vững hơn Lmean. Short windows đặc biệt có ý nghĩa khi đối tượng nghiên cứu là một quá trình chuyển tiếp. Cửa sổ ngắn cải thiện temporal localization nhưng làm giảm số state-space points, recurrence structures và trajectory pairs, từ đó tăng finite-sample uncertainty. Vì vậy, 60~s nên được xem là một practical compromise cho dataset hiện tại giữa temporal resolution và estimator stability, chứ không phải một universally optimal window. Về mặt ứng dụng, các state-space properties của short-window PPG có thể cung cấp thông tin bổ sung cho conventional HR, HRV hoặc morphology features. Tuy nhiên, nghiên cứu này không xây dựng classifier và không chứng minh real-time hay early drowsiness detection. Do đó, các descriptors như forecastability, recurrence organization và local divergence nên được xem là nền tảng cho future nonlinear feature development, chứ chưa phải bằng chứng về hiệu quả của một hệ thống phát hiện.

\subsection{Limitations}

Nghiên cứu còn một số hạn chế về dữ liệu và physiological interpretation. Cohort còn nhỏ và chỉ gồm healthy young adults, vì vậy khả năng generalize sang các population khác còn hạn chế. Awake và Drowsy được xác định từ behavioral và subjective evidence thay vì PSG, nên Drowsy không thể đồng nhất với N1 hay NREM. Ngoài ra, PPG chỉ là một peripheral cardiovascular measurement, do đó các dynamical changes và symbolic patterns không thể được xem là direct measurements của sympathetic hoặc parasympathetic activity.

Các nonlinear estimates cũng phụ thuộc vào preprocessing, phase-space reconstruction, recurrence threshold, neighborhood definition, fitting interval, data length và signal quality \cite{martingonzalez2018,sviridova2022,yeragani2004,hu2009}. Nghiên cứu hiện tại đã giảm ảnh hưởng của các yếu tố này bằng frozen parameters, quality control, parameter verification và robustness analysis, nhưng không thể loại bỏ hoàn toàn finite-data và parameter dependence. Đặc biệt, các cửa sổ 30~s vẫn cho uncertainty lớn hơn.

Cuối cùng, các findings hiện tại là associations chứ không xác lập causal physiological mechanisms. Nghiên cứu chưa đồng thời đo EEG, respiration, continuous blood pressure hay direct autonomic activity, và cũng chưa có external validation trên independent dataset. Future multimodal studies và replication trên population, sensor và protocol khác sẽ cần thiết để xác định generality của các dynamical patterns hiện tại.

Tổng thể, Awake--Drowsy transition đi kèm với những thay đổi phối hợp trong PPG forecastability, recurrence organization và local trajectory divergence. Khi đặt trong bối cảnh wake-to-sleep transition, pattern này phù hợp với một quá trình tái tổ chức của peripheral cardiovascular dynamics trong một physiological regime đang chuyển tiếp và chưa ổn định, thay vì một thay đổi đơn giản theo một trục chaos, complexity hay autonomic balance.